We are planning a shared hyperparameter-tuning system for the training
workloads in our workspace: reinforcement learning through RLlib,
Hamiltonian neural networks, neural-transport samplers, and
multi-objective models. This document surveys the tuning literature
through August 2026 at a level meant to teach, with the load-bearing
mechanisms derived in place and illustrated on one running example
rather than quoted from their papers, audits the open-source tooling
library by library against primary sources, and turns both into an
implementation plan.

The architectural conclusion is to keep Ray Tune
as the execution layer (its distributed trials, checkpointing, and
population-based training remain the strongest available, and we already
operate Ray) while owning the search layer ourselves, since Tune's
algorithm side has been in maintenance while method development moved to
Optuna, BoTorch, and the AutoML groups' libraries.

The algorithmic
conclusion is a tiered roadmap: asynchronous successive halving plus a
carefully configured Bayesian optimizer as the spine; noise-aware
multi-objective acquisition with multi-objective early stopping, expert
priors, cost-aware search, and the PB2-to-BG-PBT population line for
reinforcement learning as the differentiators; freeze-thaw scheduling
with a pre-trained learning-curve surrogate as the frontier bet. Each
recommendation is tied to the controlled evidence that supports it, and
a validation suite is specified before any of it is built.
